% SEGMENT OLP-0347-S01
% Part: lambda-calculus
% Chapter: representability
% Section: currying

\documentclass[../../../include/open-logic-section]{subfiles}

\begin{document}

\olfileid{lam}{rep}{cur}
\olsection{கரியாக்கம்}

ஒரு $\lambd$-அருவமான $\lambd[x][M]$ ஓரிடச் சார்பைப் பிரதிநிதித்துவப்படுத்துகிறது.
பல உள்ளீடுகளை ஏற்கும் சார்புகளை வரையறுக்க விரும்பும்போது இது பெரும் கட்டுப்பாடு.
இதைச் செய்வதற்கான ஒரு வழி, சோடிகள், மும்மைகள் முதலியவற்றை அமைக்க அனுமதிக்குமாறு
$\lambd$-கலனத்தை விரிவாக்குவது. அப்படியானால், எடுத்துக்காட்டாக, மூவிடச் சார்பு
$\lambd[x][M]$ தனது உள்ளீடு ஒரு மும்மையாக இருக்கும் என்று எதிர்பார்க்கும்.
ஆனால் இதை \emph{கரியாக்கம்} மூலம் செய்வது மேலும் வசதியானது.

ஓர் எடுத்துக்காட்டைக் கருதுவோம். $\lambd$-கலனத்தில் $+$ என்ற செயல் நம்மிடம்
இருப்பதாகச் சிறிது நேரம் பாவிப்போம். கூட்டல் சார்பு $2$-இடச் சார்பு; அதாவது,
அது இரண்டு உள்ளீடுகளை ஏற்கிறது. ஆனால் ஒரு $\lambd$-அருவம் ஓரிடச் சார்புகளை
மட்டுமே தருகிறது: $\lambd[(x,y)][(x+y)]$ போன்ற கோவைகளைத் தொடரியல் அனுமதிப்பதில்லை.
எனினும், $\lambd[y][(x+y)]$ ஆல் தரப்படும் ஓரிடச் சார்பு~$f_x(y)$ ஐக் கருதலாம்;
அது தனது ஒரே உள்ளீடு~$y$ உடன் $x$ ஐக் கூட்டுகிறது. உண்மையில் இது ஒரே சார்பு
அல்ல; ஒவ்வோர் எண்~$x$ க்கும் ஒன்றாக, ``$x$ ஐக் கூட்டு'' என்ற வெவ்வேறு
சார்புகளின் குடும்பம். இப்போது மற்றோர் ஓரிடச் சார்பு~$g$ ஐ $\lambd[x][f_x]$
என வரையறுக்கலாம். உள்ளீடு $x$ மீது செயல்படுத்தும்போது, $g(x)$ சார்பு~$f_x$ ஐத்
திருப்பித்தருகிறது---ஆகவே அதன் மதிப்புகளே பிற சார்புகள். இப்போது $g$ ஐ~$x$
மீதும், பின்னர் அதன் விளைவை~$y$ மீதும் செயல்படுத்தினால்,
$(g(x))y = f_x(y) = x+y$ கிடைக்கிறது. இவ்வாறு ஓரிடச் சார்பு~$g$, இருவிடக்
கூட்டல் சார்பின் அதே பணியைச் செய்ய முடியும். இருவிடச் சார்புகளை, அவற்றின்
மதிப்புகள் ஓரிடச் சார்புகளாக அமையும் ஓரிடச் சார்புகளாக மாற்றும் இந்தத்
தந்திரத்தையே ``கரியாக்கம்'' குறிக்கிறது.

$\lambd$-கலனத்தின் தொடரியலிலேயே அமைந்த ஓர் எடுத்துக்காட்டு இதோ. சார்பு
$f(x,y) = x$ ஐ எவ்வாறு பிரதிநிதித்துவப்படுத்துவது? இரண்டு உள்ளீடுகளை ஏற்று
முதலாவதைத் திருப்பித்தரும் சார்பை வரையறுக்க வேண்டுமானால்,
$\lambd[x][\lambd[y][x]]$ என்று எழுதலாம். சொற்பொருளாக இது ஓர் உள்ளீடு~$x$ ஐ
ஏற்றுச் சார்பு~$\lambd[y][x]$ ஐத் திருப்பித்தரும் சார்பு. சார்பு
$\lambd[y][x]$ மற்றோர் உள்ளீடு~$y$ ஐ ஏற்றாலும் அதை விட்டுவிட்டு எப்போதும்
~$x$ ஐத் திருப்பித்தருகிறது. $\lambd[x][\lambd[y][x]]$ ஐ இரண்டு உள்ளீடுகள்
மீது செயல்படுத்தும்போது என்ன நிகழ்கிறது என்று பார்ப்போம்:
\begin{align*}
  (\lambd[x][\lambd[y][x]])MN
  \bredone & (\lambd[y][M])N \\
  \bredone & M
\end{align*}

பொதுவாக, ஏதோ சொல்~$N$ ஆல் வரையறுக்கப்படும் $x_1$, \dots,~$x_n$ என்ற
அளவுருக்களைக் கொண்ட சார்பை எழுத,
$\lambd[x_1][\lambd[x_2][\ldots\lambd[x_n][N]]]$ ஐப் பயன்படுத்தலாம். அதன்
மீது $n$ உள்ளீடுகளைச் செயல்படுத்தினால் பின்வருவது கிடைக்கிறது:
\begin{multline*}
  (\lambd[x_1][\lambd[x_2][\ldots\lambd[x_n][N]]]) M_1 \dots M_n \bredone\\
  \begin{aligned}
  \bredone {} & (\Subst{(\lambd[x_2][\ldots\lambd[x_n][N]])}{M_1}{x_1}) M_2
  \dots M_n\\
   \eqs {} & (\lambd[x_2][\ldots\lambd[x_n][\Subst{N}{M_1}{x_1}]]) M_2
            \dots M_n \\
  \vdots & \\
  \bredone {} &\Subst{\Subst{N}{M_1}{x_1}\ldots}{M_n}{x_n}
  \end{aligned}
\end{multline*}
% SOURCE CORRECTION TA-LAM-003: the iterated substitution must begin with the body N, not an otherwise unused P.
கடைசி வரி, சார்பு வரையறையின் உடலில் $x_i$ க்குப் பதிலாக $M_i$ ஐ
இடைநிறுத்துவதைச் சொற்பொருளாகக் குறிக்கிறது; ஒரு சார்பை பல உள்ளீடுகள் மீது
செயல்படுத்தும்போது நாம் வேண்டுவதும் துல்லியமாக இதுவே.
\end{document}
